% GENERATED FILE. Edit canonical lessons, then run npm run generate:books.
% canonical-source-hash: fnv1a64:be27c460975a0071
% canonical-lessons: GE-C444-baumarkt, GE-C444-schreibwarengeschaeft, GE-C444-buchhandlung, GE-C444-blumenladen, GE-C444-kleben

\chapter{Noch mehr Läden}
\label{ch:ge-noch-mehr-laden}
\hlchaptermodality{444}

\begin{chapteropening}
\textbf{By the end of this chapter:} \emph{I can find other shops: the DIY store, the stationer, the bookshop, the flower shop and gluing things.}

Complete the last lesson of chapter 444: I can find other shops: the DIY store, the stationer, the bookshop, the flower shop and gluing things.
\end{chapteropening}

\section[der Baumarkt]{der Baumarkt — the DIY store, the hardware store}

\label{lesson:GE-C444-baumarkt}



\begin{quote}
Before the new one: say the German for the chemist, then the German for the health food shop.
\end{quote}

\subsection*{You'll want to know: der Baumarkt}
\textbf{der Baumarkt} — \textquotedblleft{}the DIY store, the hardware store\textquotedblright{}.

\begin{quote}
Im Baumarkt kaufen wir Farbe für das Bad.
\end{quote}

\begin{quote}
\emph{We buy paint for the bathroom at the DIY store.}
\end{quote}

Plural: \emph{die Baumärkte}.

\begin{cousinweb}
\emph{Baumarkt} is built from \emph{Bau} 'building', from \emph{bauen}, inherited from Proto-Germanic \emph{būaną} (EtymDB), and \emph{Markt}, from Latin \emph{mercātus} (EtymDB).
\end{cousinweb}

\begin{grammarlens}[title={What you've built}]
One more word for telling, explaining and describing at B1.
\end{grammarlens}

\subsection*{Your turn}
\begin{itemize}
\raggedright
  \item \emph{Say it:} \emph{der Baumarkt}
  \item \emph{Say it:} \emph{der Baumarkt}, once more
  \item \emph{Say it:} a sentence of your own with \emph{der Baumarkt}
  \item \emph{Recall:} say the German for the chemist, then the German for the health food shop, then say \emph{der Baumarkt} again
\end{itemize}

\begin{tcolorbox}[breakable,colback=teal!4,colframe=teal!35!black,title={Before you move on}]
What does der Baumarkt mean? (\textquotedblleft{}The DIY store, the hardware store\textquotedblright{}.) Say it once more.
\end{tcolorbox}
\section[das Schreibwarengeschäft]{das Schreibwarengeschäft — the stationer, the stationery shop}

\label{lesson:GE-C444-schreibwarengeschaeft}



\begin{quote}
Before the new one: say the German for the health food shop, then the German for the DIY store.
\end{quote}

\subsection*{You'll want to know: das Schreibwarengeschäft}
\textbf{das Schreibwarengeschäft} — \textquotedblleft{}the stationer, the stationery shop\textquotedblright{}.

\begin{quote}
Im Schreibwarengeschäft kauft das Kind neue Hefte.
\end{quote}

\begin{quote}
\emph{At the stationer the child buys new exercise books.}
\end{quote}

Plural: \emph{die Schreibwarengeschäfte}. \emph{Die Schreibwaren} — paper, pens and so on.

\begin{cousinweb}
\emph{Schreibwarengeschäft} is built from \emph{schreiben} 'to write', an early borrowing from Latin \emph{scrībere} (Lewis \& Short), \emph{Ware} 'goods', from Middle Low German \emph{ware} (EtymDB), and \emph{Geschäft} 'shop'.
\end{cousinweb}

\begin{grammarlens}[title={What you've built}]
One more word for telling, explaining and describing at B1.
\end{grammarlens}

\subsection*{Your turn}
\begin{itemize}
\raggedright
  \item \emph{Say it:} \emph{das Schreibwarengeschäft}
  \item \emph{Say it:} \emph{das Schreibwarengeschäft}, once more
  \item \emph{Say it:} a sentence of your own with \emph{das Schreibwarengeschäft}
  \item \emph{Recall:} say the German for the health food shop, then the German for the DIY store, then say \emph{das Schreibwarengeschäft} again
\end{itemize}

\begin{tcolorbox}[breakable,colback=teal!4,colframe=teal!35!black,title={Before you move on}]
What does das Schreibwarengeschäft mean? (\textquotedblleft{}The stationer, the stationery shop\textquotedblright{}.) Say it once more.
\end{tcolorbox}
\section[die Buchhandlung]{die Buchhandlung — the bookshop}

\label{lesson:GE-C444-buchhandlung}



\begin{quote}
Before the new one: say the German for the DIY store, then the German for the stationer.
\end{quote}

\subsection*{You'll want to know: die Buchhandlung}
\textbf{die Buchhandlung} — \textquotedblleft{}the bookshop\textquotedblright{}.

\begin{quote}
In der Buchhandlung findet sie einen Reiseführer.
\end{quote}

\begin{quote}
\emph{In the bookshop she finds a travel guide.}
\end{quote}

Plural: \emph{die Buchhandlungen}.

\begin{cousinweb}
\emph{Buchhandlung} is built from \emph{Buch} 'book', inherited from Old High German \emph{buoh} (EtymDB), and \emph{Handlung} in the old sense 'trade, shop', from \emph{handeln} 'to trade'.
\end{cousinweb}

\begin{grammarlens}[title={What you've built}]
One more word for telling, explaining and describing at B1.
\end{grammarlens}

\subsection*{Your turn}
\begin{itemize}
\raggedright
  \item \emph{Say it:} \emph{die Buchhandlung}
  \item \emph{Say it:} \emph{die Buchhandlung}, once more
  \item \emph{Say it:} a sentence of your own with \emph{die Buchhandlung}
  \item \emph{Recall:} say the German for the DIY store, then the German for the stationer, then say \emph{die Buchhandlung} again
\end{itemize}

\begin{tcolorbox}[breakable,colback=teal!4,colframe=teal!35!black,title={Before you move on}]
What does die Buchhandlung mean? (\textquotedblleft{}The bookshop\textquotedblright{}.) Say it once more.
\end{tcolorbox}
\section[der Blumenladen]{der Blumenladen — the flower shop, the florist's}

\label{lesson:GE-C444-blumenladen}



\begin{quote}
Before the new one: say the German for the stationer, then the German for the bookshop.
\end{quote}

\subsection*{You'll want to know: der Blumenladen}
\textbf{der Blumenladen} — \textquotedblleft{}the flower shop, the florist's\textquotedblright{}.

\begin{quote}
Für den Geburtstag kauft er Rosen im Blumenladen.
\end{quote}

\begin{quote}
\emph{He buys roses at the flower shop for the birthday.}
\end{quote}

Plural: \emph{die Blumenläden}.

\begin{cousinweb}
\emph{Blumenladen} is built from \emph{Blume} 'flower', inherited from Old High German \emph{bluomo} (EtymDB), and \emph{Laden} 'shop', first 'plank, shutter', from Middle High German \emph{laden} (EtymDB).
\end{cousinweb}

\begin{grammarlens}[title={What you've built}]
One more word for telling, explaining and describing at B1.
\end{grammarlens}

\subsection*{Your turn}
\begin{itemize}
\raggedright
  \item \emph{Say it:} \emph{der Blumenladen}
  \item \emph{Say it:} \emph{der Blumenladen}, once more
  \item \emph{Say it:} a sentence of your own with \emph{der Blumenladen}
  \item \emph{Recall:} say the German for the stationer, then the German for the bookshop, then say \emph{der Blumenladen} again
\end{itemize}

\begin{tcolorbox}[breakable,colback=teal!4,colframe=teal!35!black,title={Before you move on}]
What does der Blumenladen mean? (\textquotedblleft{}The flower shop, the florist's\textquotedblright{}.) Say it once more.
\end{tcolorbox}
\section[kleben]{kleben — to stick, to glue}

\label{lesson:GE-C444-kleben}



\begin{quote}
Before the new one: say the German for the bookshop, then the German for the flower shop.
\end{quote}

\subsection*{You'll want to know: kleben}
\textbf{kleben} — \textquotedblleft{}to stick, to glue\textquotedblright{}.

\begin{quote}
Sie klebt das kaputte Bild wieder zusammen.
\end{quote}

\begin{quote}
\emph{She glues the broken picture back together.}
\end{quote}

\emph{Der Kleber} — the glue. \emph{Das klebt} — that is sticky.

\begin{cousinweb}
\emph{kleben} is inherited from Old High German \emph{klebēn} and Proto-Germanic \emph{klibjaną} 'to stick' (EtymDB). English \emph{cleave} 'to cling' is a cousin.
\end{cousinweb}

\begin{grammarlens}[title={What you've built}]
One more word for telling, explaining and describing at B1.
\end{grammarlens}

\subsection*{Your turn}
\begin{itemize}
\raggedright
  \item \emph{Say it:} \emph{kleben}
  \item \emph{Say it:} \emph{kleben}, once more
  \item \emph{Say it:} a sentence of your own with \emph{kleben}
  \item \emph{Recall:} say the German for the bookshop, then the German for the flower shop, then say \emph{kleben} again
\end{itemize}

\begin{tcolorbox}[breakable,colback=teal!4,colframe=teal!35!black,title={Before you move on}]
What does kleben mean? (\textquotedblleft{}To stick, to glue\textquotedblright{}.) Say it once more.
\end{tcolorbox}
